% =====================================================================
%  Выводы по лабораторной работе.
% =====================================================================

\labpartbreak{Выводы}

\begin{quote}
\textit{Раздел будет окончательно сформулирован после анализа данных
из раздела «Анализ запросов из сети Internet до и после настройки» --
ниже приведён каркас выводов, который будет дополнен конкретными
числами и наблюдениями.}
\end{quote}

Семантически лабораторная работа должна показать разницу между
обнаружением и предотвращением: сам факт журналирования (лабораторная
работа №4.1) позволяет знать о происходящем, но никак на него не
влияет, тогда как активный ответ OSSEC превращает пассивное знание в
действие -- источник, уже опознанный как вредоносный, перестаёт
получать саму возможность продолжить попытки. \textit{[после сбора
данных: подтвердить конкретными числами -- сколько уникальных
источников было заблокировано за период наблюдения, сократилось ли
количество последующих попыток с тех же подсетей]}

С прагматической точки зрения работа должна продемонстрировать, что
внедрение IDS/IPS даже в простейшей, однохостовой конфигурации
(отдельный сетевой сенсор не требуется) ощутимо меняет то, с какими
данными сервер сталкивается в дальнейшем -- активный ответ не только
защищает, но и сам становится источником полезных данных
(\texttt{active-responses.log}), которых не существовало до его
установки. \textit{[после сбора данных: указать, оправдал ли себя
выбор именно HIDS-решения по сравнению с сетевыми альтернативами из
варианта задания, и с какими ограничениями пришлось столкнуться на
практике]}
